\chapter{The Algebra of Comparison}
\label{chap:algebra-of-comparison}

% Closed question: How can an instrument compare records and repeat the result
% without pretending the comparison is the thing measured?

\begin{chapteressay}

A balance scale has three admissible outputs. The left pan is lower, the
right pan is lower, or the beam is level. Three outputs, no more. The
balance cannot say how much heavier; it cannot produce a number. To
weigh an object, one adds standard reference weights to the right pan
until the balance reads ``level.'' The weight is then the sum of the
references. The weighing is a sequence of comparisons; the result is a
sum of standards; the standards are the calibration.

This is the chapter's primary observation. The instrument's natural
output is not a number; the natural output is a comparison. The
comparison is binary in the mathematical sense: it produces one of
three values for any pair of inputs. Numbers are constructed from
comparisons by combining standards. The numbers are downstream of the
comparisons.

The previous chapter introduced the partial order on marks. This chapter
asks what an instrument can do with a partial order. The answer is that
the instrument can compare two marks and produce a verdict; multiple
verdicts compose into structured studies; the studies form a lattice;
the lattice's quantitative output is a rank or percentile.

The chapter develops this through five mathematical structures. The
first is the binary interface itself: the three-valued output of a
comparison, formalized as the typeclass \texttt{BINARY}. The second is
partial-order arithmetic: meet and join operations that compose
comparisons into lattice structure. The third is repeatability: the
condition under which multiple trials of the same comparison can be
averaged. The fourth is the lattice: the algebraic structure of
admissible comparison compositions, distributive and complemented in
the Boolean case. The fifth is numerical study: the production of
rank-based quantitative outputs from comparisons alone.

Each structure is mathematical, not algorithmic. The binary interface
is a function with three possible outputs; the lattice arithmetic is
the algebra of meet and join; repeatability is a property of the
trial; the lattice is a category of algebraic objects; the numerical
study is a function from comparison sets to ranks. These are all
classical mathematical structures. The chapter's contribution is to
identify them as the natural arithmetic of comparison.

The chapter's ground example is the balance scale. The chapter's
second example is Gauss's survey of Hanover: repeated angle
measurements, averaged to improve precision by $\sqrt{n}$. The
averaging is admissible because the measurements are
\texttt{REPEATABLE}: the same angle, measured by the same
instrument, with care, produces the same reading up to a small
independent error. The $\sqrt{n}$ improvement is the central
limit theorem applied to instrument precision. It is one of
mathematics' most useful corollaries.

The chapter's second concrete example is the passport-control
queue. The kiosk's verdict is binary in the chapter's sense:
admit, refer, or deny. The verdict is not a measurement of the
traveler; it is a routing decision based on the documents and
the procedure. The procedure's algebra is the chapter's
lattice algebra; the system's quantitative output is the
processing rate and verdict distribution.

The chapter's closed question is:

\begin{center}
\emph{How can an instrument compare records and repeat the result
without pretending the comparison is the thing measured?}
\end{center}

The answer the chapter develops is the binary interface with
the algebraic lattice. The instrument compares two marks; the
comparison is binary; the comparison is repeatable; the
comparisons compose into a lattice; the lattice's output is a
rank. None of these operations require the instrument to
measure the marks directly. All require only the ability to
compare. The instrument is the algorithm of comparison; the
quantity being compared is the thing the comparison is about;
the verdict is information, not magnitude.

Subsequent chapters build on this. Chapter 3 introduces the
gap between recorded comparisons: what can be said about the
absence of explicit measurement in the interval between two
comparisons? Chapter 4 introduces the algebra of records:
how multiple records compose without collapsing. Chapter 5
introduces distance from order: how the finite count of
comparisons becomes an admissible notion of distance.

\end{chapteressay}

\section{The Binary Interface}
\label{se:the-binary-interface}

A balance scale is the canonical instrument of comparison. Two pans are
suspended from a beam. An object is placed in the left pan; a reference
object is placed in the right. The beam tilts. There are exactly three
admissible outputs of the beam's tilt: the left pan is lower (the left
object is heavier), the right pan is lower (the right object is heavier),
or the beam is level (the objects are balanced).

The balance does not say how much heavier. It does not produce a number.
It produces a ternary verdict: left, right, or balanced. The instrument's
interface is exactly this three-valued comparison. Asking the balance for
the weight of one object alone is asking the wrong question; the balance
is built to compare, not to weigh. To weigh, the operator must add or
remove standard reference weights until the balance reads ``balanced'';
the weight is then the sum of the references.

This is the chapter's first claim. The instrument's natural output is
not a number; it is a comparison. The comparison is binary in the
mathematical sense: it produces one of three values (less, equal,
greater) for any pair of inputs. Numbers are constructed from
comparisons; they are not given by the instrument directly.

Formally, a binary comparison interface for a set $S$ is a decidable
function $\le : S \times S \to \{<, =, >\}$ satisfying three
properties. First, reflexivity: $a \le a$ produces $=$ (equality)
for every $a$. Second, antisymmetry: if $a \le b$ produces $\le$
and $b \le a$ produces $\le$, then $a = b$. Third, transitivity:
if $a \le b$ produces $\le$ and $b \le c$ produces $\le$, then
$a \le c$ produces $\le$.

These three properties define a total order when every pair is
comparable, and a partial order when some pairs are incomparable. The
balance scale produces a total order on weights: any two physical
objects can be compared (one is heavier, or they balance). Other
instruments produce partial orders: a chemical assay may detect
that compound A contains element X and compound B contains element
Y, but cannot directly compare A and B if their element profiles
overlap.

In Lean, the binary interface is the typeclass \texttt{BINARY} in
\texttt{Episode2.lean}:

\begin{leanexcerpt}{BINARY}
class BINARY
    (Value : Type)
    (Carrier : CarrierProcess Value)
    [d : DISTINGUISHABLE Value Carrier]
    [a : ADMISSIBLE Value Carrier]
    [c : COUNTABLE Value Carrier]
    [e : ENCODED Value Carrier]
    [r : RESIDUE Value Carrier]
  where
  observation_process : ObservationProcess Value Carrier
  zero : Limit
  one  : Limit
  bit  : Sample
  different? : Sample $\to$ Sample $\to$ Prop := ...
\end{leanexcerpt}

\texttt{BINARY} sits on top of the prior cascade (DISTINGUISHABLE,
ADMISSIBLE, COUNTABLE, ENCODED, RESIDUE). It carries the
\texttt{ObservationProcess}, two reference points \texttt{zero}
and \texttt{one}, the current \texttt{bit}, and the predicate
\texttt{different?} that decides whether two samples are
distinguishable. The compiler enforces the stack: \texttt{BINARY}
cannot be admitted without instances of every prior class.

The \texttt{Trial} type carries one application of the comparison:
two inputs and the resulting verdict. The \texttt{Trial.le}
proposition asserts that the result was $\le$ (the first input is
at most the second). The \texttt{Trial.lt} proposition asserts
strict less-than. The trial is the chapter's algebraic atom: the
record of one admissible comparison.

The chapter's second claim is that the comparison is not the
quantity being compared. The balance produces a verdict; the
weight is what the verdict is about. The two are distinct. The
verdict can be repeated reliably; the weight is the thing whose
existence the verdict witnesses. Conflating them is a category
error: the verdict is a piece of information about the weights;
the weights are physical properties of the objects.

\section{Partial Order Arithmetic}
\label{se:partial-order-arithmetic}

The divisibility lattice on the positive integers carries enough
arithmetic to support the chapter's algebra. The meet of two elements
$a, b$ is the greatest common divisor: $a \wedge b = \gcd(a, b)$. The
join is the least common multiple: $a \vee b = \operatorname{lcm}(a, b)$.

For $a = 12$ and $b = 18$: the divisors of 12 are $\{1, 2, 3, 4, 6,
12\}$; the divisors of 18 are $\{1, 2, 3, 6, 9, 18\}$. The common
divisors are $\{1, 2, 3, 6\}$. The greatest is $6$. So $12 \wedge 18 =
6$. The multiples of 12 are $\{12, 24, 36, 48, \ldots\}$; the
multiples of 18 are $\{18, 36, 54, 72, \ldots\}$. The common multiples
are $\{36, 72, 108, \ldots\}$. The least is $36$. So $12 \vee 18 =
36$. Note that $\gcd(12,18) \cdot \operatorname{lcm}(12,18) = 6 \cdot
36 = 216 = 12 \cdot 18$. This is a general identity:
$\gcd(a,b) \cdot \operatorname{lcm}(a,b) = a \cdot b$.

The meet and join operations together make the divisibility relation
into a \emph{lattice}: a partially ordered set in which every pair of
elements has both a least upper bound and a greatest lower bound. A
lattice may be distributive ($a \wedge (b \vee c) = (a \wedge b) \vee
(a \wedge c)$) or not. The divisibility lattice is distributive: the
prime factorizations distribute over both gcd and lcm.

In Lean, a lattice instance carries the meet and join operations
together with proofs of the algebraic laws (commutativity,
associativity, absorption, idempotence, and possibly distributivity).
The \texttt{Mathlib.Order.Lattice} module provides the typeclass
hierarchy: \texttt{Lattice}, \texttt{DistribLattice},
\texttt{BoundedLattice}, and so on. Each typeclass adds new laws to
the previous ones.

The chapter's claim is that the lattice structure is the natural
arithmetic of comparisons. The meet and join are not numerical
quantities; they are the results of admissible comparison operations.
$\gcd(12, 18)$ is the largest element that compares less than or
equal to both $12$ and $18$ under divisibility. $\operatorname{lcm}$
is the smallest element that compares greater than or equal to both.
The two operations are dual.

The \texttt{Study} inductive in \texttt{Episode2.lean} captures the
algebraic content of a structured inquiry:

\begin{leanexcerpt}{Study}
inductive Study
  | hypothesis : Fact $\to$ Study
  | data       : Fact $\to$ Trial $\to$ Study $\to$ Study
\end{leanexcerpt}

A \texttt{Study} is either a \texttt{hypothesis} (a Fact) or a \texttt{data}
record extending a prior study with one new \texttt{Trial}. The structure
is recursive: a study accumulates trials by adding one at a time. The
\texttt{Study.le} partial order says: one study is at most another when
the prior study fits into the other's recursive structure with matching
facts and trial ordering.

This is the chapter's algebraic version of comparison composition.
Multiple comparisons compose into a structured study; the studies
form a lattice under the implication order. The meet of two studies
is the common information; the join is their combined conclusion.
The algebra of comparisons is the algebra of studies.

A practical example. Suppose the balance scale is used to compare
five objects $A, B, C, D, E$. The trials are: $A < B$, $B < C$,
$C = D$, $D < E$. The study's conclusion is the total order $A < B
< C = D < E$. The conclusion is derived from the trials by
transitivity (a property of \texttt{BINARY} instances). The
\texttt{Study} carries both the trials and the conclusion; the
conclusion is computable from the trials, but the trials are the
record of the underlying comparisons.

If a new trial is added --- say $A = C$ --- the trials become
inconsistent (we have $A < B < C = D$ but $A = C$ would imply
$A < B < A$, a contradiction). The \texttt{Study} rejects the new
trial: the conclusion cannot be updated to include both the old
trials and the new one. The rejection is the chapter's algebraic
content of comparison admissibility.

\section{Repeatability}
\label{se:repeatability}

Carl Friedrich Gauss, in his survey of the Kingdom of Hanover from 1818
to 1826, measured each principal angle of the triangulation network many
times and averaged the readings. The reason was statistical: the
individual measurements were subject to small errors of vision, of
levelness, of atmospheric refraction. The average of many measurements
was more reliable than any single one.

The mathematical statement is precise. Suppose the true angle is
$\theta$, and each independent measurement is $\theta_i = \theta +
\epsilon_i$, where $\epsilon_i$ is a small error drawn independently
from a distribution with mean zero and standard deviation $\sigma$. The
average of $n$ measurements is $\bar{\theta} = \theta + \frac{1}{n}
\sum \epsilon_i$. The standard deviation of $\bar{\theta}$ is
$\sigma / \sqrt{n}$. Ten measurements improve precision by a factor of
$\sqrt{10} \approx 3.16$; one hundred measurements by a factor of ten.

The square-root improvement is one of mathematics' most important laws.
It is the central limit theorem applied to instrument precision. It
makes averaging admissible: combining many measurements produces a
more precise estimate than any one. The improvement is not arbitrary;
it follows from the assumption that the errors are independent.

The condition for averaging is repeatability: each measurement must be
an independent admissible trial of the same underlying quantity. If
the measurements share a systematic bias (a constant error that does
not average to zero), averaging does not remove it; only the random
component shrinks. If the measurements are correlated (an error in one
measurement makes the next more likely to have the same error),
averaging fails to converge to the true value.

Repeatability is therefore a structural condition. The instrument must
produce comparable outputs across independent trials. The
\texttt{REPEATABLE} typeclass in \texttt{Episode2.lean} formalizes
this:

\begin{leanexcerpt}{REPEATABLE}
class REPEATABLE
    (Value : Type)
    (Carrier : CarrierProcess Value)
    [d : DISTINGUISHABLE Value Carrier]
    [a : ADMISSIBLE Value Carrier]
    [c : COUNTABLE Value Carrier]
    [e : ENCODED Value Carrier]
    [r : RESIDUE Value Carrier]
    [b : BINARY Value Carrier]
  where
  repeatable_process : RepeatableProcess Value Carrier
  typical_response   : Trial $\to$ Trial $\to$ Prop := ...
\end{leanexcerpt}

\texttt{REPEATABLE} is admissible only on top of the full prior cascade
including \texttt{BINARY}. It carries the underlying
\texttt{RepeatableProcess} and the \texttt{typical\_response} predicate
that says when one trial's outcome is the expected continuation of
another. The compiler enforces the cascade: an instance is constructible
only after every prior class has been provided.

Repeatability is the condition under which the chapter's algebra of
comparisons becomes meaningful. Without it, two studies might
produce different conclusions from the same trials; the
\texttt{Study} type cannot certify reproducibility. With it, the
algebra closes: the same input always produces the same comparison,
and the same set of comparisons always produces the same study.

The chapter's claim is that repeatability is the precondition for
inference. From repeatable comparisons, the instrument can build
studies. From studies, the inferring agent can draw conclusions.
From conclusions, theories are constructed. The chain of
inference rests on the repeatability of the comparisons at the
base.

Repeatability has its limits. Quantum measurements are not
\texttt{REPEATABLE} in the typeclass's sense: two measurements of
the same observable on identically prepared states can produce
different outcomes. The probability distribution of outcomes is
reproducible (over many trials), but individual outcomes are not.
Quantum instruments produce \texttt{REPEATABLE} \emph{statistics},
not \texttt{REPEATABLE} individual readings. The chapter's
discussion of probabilistic repeatability is taken up in later
chapters.

For classical instruments --- balances, thermometers, rulers, clocks
--- the deterministic version of \texttt{REPEATABLE} is the
standard. The instruments are designed to produce the same reading
on the same input. When they fail to do so, the failure is a
calibration error or a hardware fault, not a feature of the
measurement.

Gauss's procedure is the canonical model. Each angle was measured
many times. The measurements were combined by averaging. The
averaged reading was more reliable than any individual measurement.
The procedure is admissible because the measurements were
\texttt{REPEATABLE}. The same procedure applied to a non-repeatable
instrument would produce a number whose meaning could not be
verified.

\section{The Lattice}
\label{se:the-lattice}

The power set of a set $S$ is the set of all subsets of $S$. Under
inclusion ($A \subseteq B$), the power set is a partial order. The
meet of two subsets is their intersection: $A \wedge B = A \cap B$.
The join is their union: $A \vee B = A \cup B$. The complement is
set difference: $A^c = S \setminus A$. The structure is a
\emph{Boolean algebra}: a distributive lattice with a complement
operation.

What the power-set Boolean algebra adds that the divisibility
lattice did not: \emph{complementation}. The divisibility lattice
on positive integers has gcd and lcm, but it has no complement
operation: there is no natural number $n^c$ such that
$\gcd(n, n^c) = 1$ and $\operatorname{lcm}(n, n^c)$ equals some
fixed top element. The divisibility lattice has no top element at
all (the lcm of all positive integers is not a positive integer).
The power-set lattice has both: a top element $S$ and, for every
$A$, a complement $A^c$ such that $A \cap A^c = \emptyset$ and
$A \cup A^c = S$. The complement is the operation that lets the
algebra of comparisons express \emph{negation}, which the
divisibility lattice cannot.

For $S = \{1, 2, 3\}$, the power set has eight elements: $\emptyset,
\{1\}, \{2\}, \{3\}, \{1,2\}, \{1,3\}, \{2,3\}, \{1,2,3\}$. The
Hasse diagram is a cube. The bottom is $\emptyset$; the top is $S$
itself. Each edge in the diagram is one element added or removed.
The diagram has eight vertices, twelve edges, six faces, one
interior region.

The distributive law holds: $A \cap (B \cup C) = (A \cap B) \cup (A
\cap C)$. The proof is element-chasing: an element is in the left
side iff it is in $A$ and in $B \cup C$, iff it is in $A$ and in
$B$, or in $A$ and in $C$, iff it is in the right side.

The complement law holds: $A \cap A^c = \emptyset$ and $A \cup A^c
= S$. Every element is either in $A$ or in $A^c$, never both, and
the union covers $S$.

These laws are the standard axioms of Boolean algebra. Every Boolean
algebra is isomorphic to the power set of some set (Stone's
representation theorem). The chapter's claim is that the algebra of
comparisons is precisely this Boolean algebra: comparisons compose
via meet and join, complements correspond to negations of
propositions, and the laws ensure the algebra is consistent.

The lattice structure is not unique to subset inclusion. Many
partial orders give rise to lattices. The divisibility lattice on
positive integers is one (meet = gcd, join = lcm). The lattice of
propositions in classical logic is another (meet = conjunction,
join = disjunction, complement = negation). The lattice of
sub-modules of a module is a third. Each is a different concrete
realization of the same algebraic structure.

For comparison processes, the lattice structure means: studies can
be combined via meet (the common conclusion of two studies) and
join (the combined conclusion). The meet is the conservative
combination: only conclusions implied by both studies are
retained. The join is the additive combination: all conclusions
from both studies are kept.

When two studies are inconsistent --- their joint conclusions imply
a contradiction --- the join is the entire lattice (the trivial
``everything is true'' state), which is a signal that the
combination is inadmissible. The lattice's structure includes
this top element as a marker of failure: a study whose conclusions
include everything has lost all informational content.

In Lean, the lattice structure for studies is provided by the
\texttt{Study.le} relation in \texttt{Episode2.lean}. The relation
makes \texttt{Study} a poset. Mathlib's lattice infrastructure
extends the poset to a lattice when the meet and join exist. For
finite studies (with a finite number of trials), the meet and join
always exist and can be computed.

The chapter's claim is that the lattice structure is the natural
arithmetic of comparison. Algebraic operations on comparisons
(combine, intersect, complement) are well-defined; the algebra is
distributive; the algebra has a top and bottom element. The
arithmetic of comparison is not numerical arithmetic; it is
lattice arithmetic. The two are different algebras.

\section{Numerical Study}
\label{se:numerical-study}

The median of a finite set of numbers is the middle element after sorting.
For five numbers $\{3, 7, 8, 12, 20\}$, the median is 8 (the third
element). For four numbers $\{3, 7, 12, 20\}$, the median is conventionally
the average of the two middle elements: $(7 + 12)/2 = 9.5$. The median is
a comparison-based statistic: it depends only on the order of the inputs,
not on their numerical values.

This is distinct from the mean, which is an arithmetic statistic. The mean
of $\{3, 7, 8, 12, 20\}$ is $50/5 = 10$. The mean depends on the
numerical values; it requires arithmetic operations (addition, division).
The median depends only on the order; it requires only the ability to
compare.

The median is robust to outliers. Adding a single outlier to a dataset
(say, $\{3, 7, 8, 12, 20, 1000\}$) shifts the mean significantly (from
$10$ to about $175$) but shifts the median only modestly (from $8$ to
$10$). The robustness is because the comparison-based statistic does
not weight extreme values more than central ones; each value
contributes its position in the order, not its magnitude.

The chapter's claim is that the median is the canonical comparison-based
statistic. It is the center of the comparison order, not the center of
mass. Other comparison-based statistics include quartiles, percentiles,
and the interquartile range. All of these depend only on ranking, not
on arithmetic. They are admissible inferences from comparison data.

In Lean, the median and related statistics are part of the
\texttt{NUMERIC} typeclass cascade in \texttt{Episode2.lean}. A process
is \texttt{NUMERIC} when it carries the comparison stack and produces
an admissible study-level output. The \texttt{ComputationalProcess}
structure underlying \texttt{NUMERIC} carries the
\texttt{RepeatableProcess} and an optional output \texttt{Study}.

\begin{leanexcerpt}{ComputationalProcess}
structure ComputationalProcess
    (Value : Type)
    (Carrier : CarrierProcess Value)
    [d : DISTINGUISHABLE Value Carrier]
    [a : ADMISSIBLE Value Carrier]
    [c : COUNTABLE Value Carrier]
    [e : ENCODED Value Carrier]
    [r : RESIDUE Value Carrier]
    [b : BINARY Value Carrier]
    [e : REPEATABLE Value Carrier]
  where
  repeatable_process : RepeatableProcess Value Carrier
  output : Option Study
  error_code : Study := ...
  close   : Option Study $\to$ Study := ...
  closure : Study $\to$ Study := ...
\end{leanexcerpt}

A \texttt{ComputationalProcess} requires the full prior cascade
through \texttt{REPEATABLE}. It carries the repeatable process, an
optional \texttt{Study} output, an error fallback, and the closure
rule that turns a pending result into a final study. The chapter's
admissible numerical output is the closed \texttt{Study}.

The chapter closes here. The instrument compares; the comparisons
are repeatable; the comparisons compose into studies; the studies
form a lattice; the studies admit comparison-based statistics. The
algebraic structure of comparison is the lattice; the
quantitative output is the rank, percentile, or other ranking
statistic. None of this requires arithmetic on the underlying
values; all of it requires only the ability to compare.

This is the chapter's answer to its closed question. The instrument
compares records by producing one of three verdicts ($<$, $=$, or
$>$) for any pair of inputs. The comparisons are repeatable: the
same input pair produces the same verdict each time. Multiple
comparisons compose into structured studies via lattice operations
(meet, join). The numerical output of a study is a rank or
percentile, computed from comparisons alone. The verdict is not the
quantity being compared; the verdict is the algebraic record of the
comparison.

The instrument is the algorithm that produces the verdict. The
quantity being compared is the thing the verdict is about. The
algebraic structure of comparison is the lattice of studies; the
quantitative output is the rank. The chapter has not produced any
direct numerical measurement of the original quantities. It has
produced an admissible algebra of comparisons. Subsequent chapters
will build on this algebra: distance from order, motion from
selection, transport between studies, and the calibration that
makes multiple studies commensurable.

\begin{bridge}

The chapter's question was how an instrument can compare records and
repeat the result without pretending the comparison is the thing
measured.

The answer is the binary interface. The instrument produces a
ternary verdict ($<$, $=$, $>$) for any pair of inputs. The verdict
is reproducible: the same input pair produces the same verdict. The
verdicts compose via the partial-order algebra of the lattice; the
lattice operations (meet, join) are the algebraic structure of
comparison. Repeatability is the precondition: only reproducible
verdicts can be averaged or combined. The numerical output of a
collection of verdicts is a rank or percentile, computed from
comparisons alone. The verdict is not the quantity being compared;
the verdict is information about the quantity, expressed in the
language of the comparison.

What remains is the gap between recorded events. The chapter has
recorded a sequence of comparisons, but between any two
comparisons, the instrument carries no record. What can be said
about that gap? Chapter 3 takes up the question: what is the
admissible structure between recorded events, given only the marks
and the order?

\end{bridge}

\begin{coda}{The Passport-Control Queue}

The airport terminal is busy this morning. International arrivals are
queueing at the passport-control kiosks. Each kiosk has an officer behind
a glass partition, a small electronic display showing the queue position,
and a slot for the passport. Travelers approach in line and present their
documents one at a time.

The kiosk's interface is binary in the chapter's sense. For each
traveler, the officer's decision is one of three admissible outputs:
\emph{admit} (the traveler is cleared to enter), \emph{refer to secondary
inspection} (the documents require additional verification), or \emph{deny
entry} (the traveler must leave on the next outbound flight).

The officer does not produce a number. The officer produces a ternary
verdict. The passport is examined; the database is queried; the
photograph is compared to the traveler's face; the questions are asked
and answered. The result is one of three outcomes. The kiosk's
display registers the outcome; the queue advances; the next traveler
approaches.

This is the binary interface in institutional form. The officer's
inputs are the documents and the traveler. The output is the ternary
verdict. The verdict is not a measurement of the traveler's
citizenship, identity, or character; the verdict is the routing
decision based on the admissible documents and the answers to the
admissible questions. The traveler is not in any sense ``measured'' by
the officer. The traveler is routed.

The decision is repeatable in the chapter's strict sense. The same
traveler with the same documents, examined by the same officer (or by
any officer trained in the same procedures), should produce the same
verdict. If two officers produce different verdicts for the same
documents, the system has a calibration problem; the procedures must
be reviewed. The reproducibility is the institutional version of the
\texttt{REPEATABLE} typeclass.

The algebra of the kiosk's verdicts is a small lattice. The set of
all admissible verdicts is $\{admit, refer, deny\}$, ordered: admit
$\le$ refer $\le$ deny (the verdicts are arranged by their
implications for the traveler). The meet of two verdicts is the
more lenient: if officer A says admit and officer B says refer, the
joint verdict is the more lenient (admit), unless an overruling
procedure intervenes. The join is the stricter: if A says admit and
B says deny, the joint verdict is the stricter (deny).

In practice the kiosk does not perform meets and joins of multiple
verdicts; each traveler is processed by one officer. But the
algebraic structure is the abstract framework within which the
procedures are designed. The system architects choose, in advance,
which combinations of evidence produce which verdicts. The choices
are the calibration of the procedure.

A traveler appears at the kiosk with an electronic visa that has been
issued by the destination country's consulate. The officer queries the
database. The visa is valid. The traveler's photograph matches. The
verdict is \emph{admit}. The traveler proceeds through the gate to
the baggage claim area.

A traveler appears with an expired visa. The officer queries the
database. The visa expired four months ago. The verdict is \emph{deny
entry}. The traveler is escorted to the airline's representative, who
arranges the next outbound flight. The denial is admissible; the
underlying physical fact is that the visa expired before the
arrival. The verdict is information about the visa, not about the
traveler in any deeper sense.

A traveler appears with an unusual travel pattern: the passport shows
recent visits to several countries known to be conflict zones. The
officer's training has prepared for this case: the verdict is
\emph{refer to secondary inspection}. The traveler is taken to a
separate room where additional questioning occurs. The secondary
officer has the same binary interface but with a more thorough
procedure. The eventual verdict may be admit or deny.

The chapter's claim is that the algebra of the kiosk is the
algebra of comparisons made admissible by procedure. The officer's
training is the procedure; the procedure is the calibration; the
verdicts are the admissible outputs. The procedure is
\texttt{REPEATABLE}: two officers applying the same procedure to
the same case should reach the same verdict.

The system's quantitative output --- the numerical study --- is the
processing rate. The terminal handles roughly five hundred
travelers per hour; the kiosks together handle that flow. The
processing rate is the rank of the system: how many travelers per
hour, on average, are processed at each verdict tier. The verdict
distribution (what percentage of travelers are admitted, referred,
or denied) is the system's percentile output.

These quantitative outputs are admissible because they are computed
from comparisons (admit/refer/deny) without requiring numerical
inputs from the travelers themselves. The travelers do not provide
a number; they provide documents. The documents are compared to
the database; the comparisons produce verdicts; the verdicts are
counted into the system's rate statistics.

At noon, the morning rush ends. The line shortens. The officers
process travelers more leisurely. The percentile statistics for
the morning are tallied: 92 percent admitted, 7 percent referred,
1 percent denied. The denial rate is within the airport's expected
range; the referral rate is slightly elevated due to a recent
travel advisory; the admission rate is normal.

The supervisor reviews the statistics at the end of the shift. The
review is itself a comparison: this morning's rates versus the
average for similar days. The comparison admits a verdict: normal,
anomalous, or requiring further investigation. The supervisor's
verdict for this morning is ``normal.'' The shift closes.

The passport-control queue performs the chapter's structure in
operational form. The binary interface is the kiosk's verdict
system. Repeatability is the officers' training. The lattice
algebra is the framework for combining evidence. The numerical
study is the processing rate and verdict distribution. None of the
operations require the officer to measure the traveler; all of the
operations require admissible comparisons against the procedure's
calibration.

The chapter's distinction --- that the comparison is not the
quantity being compared --- is what the kiosk's procedure
embodies. The officer does not say what the traveler is. The
officer says what the documents allow. The verdict is information
about the documents, expressed in the language of the procedure.
The traveler is the thing the documents are about; the procedure
is the algebra of admissible comparisons.

\end{coda}
